\chapter{EVALUATION AND REFLECTION}

\section{Evaluation of the Internship Program}

\placeholder{Evaluate the program's relevance, structure, mentorship,
technical depth, learning value, and connection to the student's major.}

\section{Impression of the Company and Work Environment}

\placeholder{Describe the professional environment using fair and verified
observations. Avoid unsupported marketing claims.}

\section{Strengths and Areas for Improvement}

\begin{table}[H]
  \centering
  \caption{Qualitative self-evaluation}
  \begin{tabularx}{\textwidth}{@{}P{4.0cm}Y@{}}
    \toprule
    \textbf{Area} & \textbf{Assessment} \\
    \midrule
    Strengths & \placeholder{Specific strengths supported by evidence} \\
    Operational discipline & \placeholder{Planning, testing, documentation,
      communication, or handover practices} \\
    Problem solving & \placeholder{Research, diagnosis, and solution approach} \\
    Technical communication & \placeholder{Writing, presentation, or
      explanation skills} \\
    Areas for improvement & \placeholder{Honest, actionable improvement areas} \\
    \bottomrule
  \end{tabularx}
\end{table}

\section{Recommendations}

\placeholder{Provide constructive recommendations for future students, the
internship program, the company, or the student's own future work.}